\section{Continuous Action Spaces}
% The old opening slide ("Can reinforcement learning solve robotics?") pasted
% three cropped paper excerpts from AlphaGo Zero / OpenAI Five / AlphaStar just
% to say "their action spaces are large but discrete". That single point is made
% directly, in one bullet, on the slide below --- so the excerpt slide is gone.

\begin{frame}{What is \emph{continuous control}?}
\vspace{0.6em}
\textbf{Control:} choosing actions over time to drive a dynamical system toward desired behaviour --- exactly the RL problem on an MDP $(\mathcal{S},\mathcal{A},P,r,\gamma)$ from Day 1.
\vspace{0.6em}
\begin{itemize}
    \item \textbf{Discrete control} (Days 1--2): the action set $\mathcal{A}$ is a \emph{finite list} --- move left/right, an Atari button, one of Go's 4{,}672 legal moves. The policy picks one entry.
    \vspace{0.3em}
    \item \textbf{Continuous control} (today): the action is a \emph{real-valued vector} $a \in \mathbb{R}^{d}$, usually bounded, e.g.\ $\mathcal{A} = [-1,1]^{d}$.
    \begin{itemize}
        \footnotesize
        \item Robot arm: joint torques. \quad Locomotion: motor commands. \quad Driving: steering angle + throttle.
        \item The set is \emph{uncountable} --- there is no list to enumerate.
    \end{itemize}
    \pause
    \vspace{0.5em}
    \item \textbf{Why this is not just a bigger discrete problem.} Go, Dota and StarCraft have \emph{enormous but still discrete} action sets ($10^{26}$ choices per step in StarCraft) --- you can in principle enumerate them and take an $\arg\max$. In continuous control you cannot enumerate at all.
\end{itemize}
\vspace{0.4em}
\centering\footnotesize\textcolor{red}{This lecture: how to do RL when actions are continuous.}
\end{frame}
